\begin{tikzpicture}[
  font=\scriptsize,
  bagian/.style={draw, rounded corners, fill=praditagreen!15, text width=2.3cm,
                 minimum height=1.3cm, align=center, font=\scriptsize\bfseries},
  tonggak/.style={draw, rounded corners, fill=blue!8, text width=2.3cm,
                  minimum height=2.3cm, align=left, font=\scriptsize},
  panah/.style={-{Stealth[length=2mm]}, thick}
]
  \node[bagian] (b1) at (0,0)
    {Bagian I\\Strategi platform\\{\tiny\mdseries Pertemuan 1 sampai 3}};
  \node[bagian] (b2) at (3,0)
    {Bagian II\\Kondisi saat ini\\{\tiny\mdseries Pertemuan 4 sampai 7}};
  \node[bagian] (b3) at (6,0)
    {Bagian III\\Platform API target\\{\tiny\mdseries Pertemuan 8 sampai 11}};
  \node[bagian] (b4) at (9,0)
    {Bagian IV\\Nilai dan tata kelola\\{\tiny\mdseries Pertemuan 12 sampai 14}};

  \node[tonggak, fill=orange!12] (c1) at (0,-2.35)
    {\textbf{C1 Konteks dan Strategi}\\[2pt]
     {\tiny profil organisasi, daftar pertukaran data, kanvas platform, peta
     ekosistem mitra}};
  \node[tonggak] (c2) at (3,-2.35)
    {\textbf{C2 Kondisi Saat Ini}\\[2pt]
     {\tiny inventaris integrasi, peta integrasi, risiko keamanan API, kesiapan
     standar sektor}};
  \node[tonggak] (c3) at (6,-2.35)
    {\textbf{C3 Kondisi Target}\\[2pt]
     {\tiny portofolio API, kontrak OpenAPI, manajemen API, portal pengembang}};
  \node[tonggak, fill=green!10] (c4) at (9,-2.35)
    {\textbf{C4 Peta Jalan dan Tata Kelola}\\[2pt]
     {\tiny model monetisasi, tata kelola API, peta jalan platform}};

  \draw[panah] (b1) -- (b2);
  \draw[panah] (b2) -- (b3);
  \draw[panah] (b3) -- (b4);
  \foreach \nomor in {1,2,3,4} {
    \draw[panah] (b\nomor) -- (c\nomor);
  }

  \node[font=\tiny, orange!60!black, anchor=north] at (c1.south)
    {bab ini dimulai di sini};
  \node[font=\tiny, gray, anchor=north] at (c2.south)
    {ditinjau menjelang UTS};
  \node[font=\tiny, gray, anchor=north] at (c4.south)
    {dipresentasikan menjelang UAS};
\end{tikzpicture}
